\section{Introducción}\label{sec:intro}

La operación diaria y la programación subestacional de sistemas eléctricos hidrotérmicos con participación creciente de generación solar fotovoltaica (FV) requieren pronósticos determinísticos confiables desde el horizonte de un día ($D+1$) hasta cuatro semanas ($7$ a $28$ días). En el Sistema Interconectado Nacional (SIN) de Colombia, este problema plantea dos exigencias simultáneas: anticipar la producción de centrales individuales con historial propio y estimar la generación de emplazamientos o zonas geográficas nuevas a partir de plantas de referencia. En ambos casos, la capacidad nominal instalada resume el tamaño de la planta bajo condiciones estándar de ensayo, pero resulta insuficiente para predecir la energía entregada, pues la irradiancia incidente, la temperatura de celda, la nubosidad convectiva tropical, la relación DC/AC, las ampliaciones escalonadas de potencia neta y las restricciones operativas separan la producción observada de la capacidad teórica de placa \cite{mueller2016yield,trigo2021models,purohit2018performance}.

Cuando se modela una flota heterogénea en lugar de una sola instalación, la formulación directa en unidades de energía bruta ($\text{kWh/día}$) introduce una distorsión estructural inmediata. En la cohorte analizada en este estudio ---$17$ plantas reales del mercado eléctrico colombiano entre enero de 2024 y octubre de 2026--- la capacidad efectiva neta varía entre $9{,}8$~MW y $388{,}0$~MW, lo que representa un cociente de escala de $39{,}6\times$. Bajo una función de pérdida calculada en $\text{kWh}$ sin normalización previa, los residuos de una central de gran escala dominan el gradiente frente a los de plantas pequeñas, impidiendo transferir parámetros entre zonas geográficas. Al mismo tiempo, los registros operativos reales contienen colas asimétricas hacia cero ocasionadas por días incompletos, desconexiones parciales, nubosidad severa o restricciones operativas, frente a los cuales la pérdida cuadrática ordinaria resulta excesivamente sensible, y los pronósticos numéricos del tiempo (NWP) pierden habilidad predictiva conforme el horizonte se extiende hacia dos, tres y cuatro semanas.

Frente a estas dificultades, este trabajo formaliza, analiza y establece el protocolo de evaluación del modelo \textit{Structured Monotone Capacity-Factor Huber Regressor} (\textbf{SM-CFHR}). El estimador responde a una secuencia de decisiones constructivas motivadas por los datos: normaliza tanto los niveles históricos recientes como la variable objetivo al espacio adimensional del factor de capacidad mediante la energía nominal máxima vigente al origen de predicción, $24\,000\,C_{i,t}$; estabiliza la respuesta frente a episodios anómalos estandarizándola con mediana y desviación absoluta mediana (MAD) estimadas exclusivamente en entrenamiento, y minimiza una pérdida de Huber \cite{huber1964,huber2009robust}; introduce una penalización cuadrática diferenciada entre dinámica temporal/calendario ($\alpha_{\mathrm{lin}}$) y variables meteorológicas ($\alpha_{\mathrm{meteo}}$); e impone restricciones de no negatividad sobre los coeficientes de persistencia de treinta días y de irradiación horizontal pronosticada, resolviendo el programa convexo resultante mediante L-BFGS-B con gradiente analítico exacto \cite{byrd1995lbfgsb,zhu1997lbfgsb}.

Dado que la arquitectura de SM-CFHR se fijó tras inspeccionar el comportamiento del panel histórico 2024--2026, la evaluación fuera de muestra (OOF) sobre ese mismo periodo debe interpretarse bajo la advertencia de reutilización adaptativa \cite{cawley2010overfit} y no equivale a una muestra confirmatoria prospectiva independiente.

Las contribuciones verificables del trabajo son tres: (1) la formulación de SM-CFHR basada en normalización causal por capacidad vigente ($24\,000\,C_{i,t}$), regularización anisotrópica y restricciones de monotonía integrada en un protocolo multihorizonte sin fuga; (2) el análisis convexo y verificación formal en Lean~4 de sus propiedades sobre $\mathbb{R}$; y (3) la evaluación fuera de muestra en $28$ configuraciones experimentales frente a $22$ referencias en energía bruta, junto con la reconciliación inferencial sobre $\mathrm{MAE}_{\mathrm{macro}}$ y $\mathrm{MAE}_{\mathrm{CF,macro}}$ que delimita la ventaja confirmada al régimen espaciotemporal con meteorología y documenta explícitamente sus derrotas en estrategia directa intraplanta sin clima.
